\section{State-complexity consequences}\label{sec:state-complexity}

Theorem~\ref{thm:FM-duality} converts trap thresholds into memory requirements:
\[
 F_{\mathcal C,\rho}(s)>n
 \quad\Longleftrightarrow\quad
 M_{\mathcal C,\rho}(n)\le s.
\]
The strict inequality on the left is essential.

\subsection{Consequences of the exact low-state thresholds}

For all induced grid mazes at radius one, Section~\ref{sec:memoryless-radius} gives $F(1)=4$, while the central theorem of Section~\ref{sec:two-state} gives $F(2)=7$.  Therefore
\begin{equation}\label{eq:M-original-small}
 \boxed{M(1)=M(2)=M(3)=1},
 \qquad
 \boxed{M(4)=M(5)=M(6)=2},
\end{equation}
and
\begin{equation}\label{eq:M7-bound}
 \boxed{3\le M(7)<\infty}.
\end{equation}
Indeed, $F(1)>n$ exactly for $n\le3$ among these horizons, and $F(2)>n$ exactly for $n\le6$.  The strict lower bound at $n=7$ follows because $F(2)=7$ is not greater than seven; finiteness follows from Theorem~\ref{thm:finite-M} below.

For induced paths, Theorem~\ref{thm:path-two-state} and Corollary~\ref{cor:path-infinity} give the sharper hierarchy
\begin{equation}\label{eq:M-path-small}
 M_{\mathcal P,1}(n)=1\quad(n\le3),
 \qquad
 M_{\mathcal P,1}(n)=2\quad(4\le n\le7),
\end{equation}
and
\begin{equation}\label{eq:M-path-large}
 \boxed{M_{\mathcal P,1}(n)\in\{3,4\}\qquad(n\ge8).}
\end{equation}
The first equality is deductive; the exact value $2$ on $4\le n\le7$ inherits the finite lower certificate of Theorem~\ref{thm:path-two-state}. The two-value alternative \eqref{eq:M-path-large} is deductive: the analytic bound $F_{\mathcal P,1}(2)\le8$ gives $M>2$, and the four-state universal controller gives $M\le4$.

\begin{table}[!htbp]
\centering\small
\begin{tabular}{@{}c|c|c@{}}
\toprule
Size horizon & all induced grid mazes, $M(n)$ & induced paths, $M_{\mathcal P,1}(n)$\\
\midrule
$1\le n\le3$ & $1$ & $1$\\
$4\le n\le6$ & $2$ & $2$\\
$n=7$ & $3\le M(7)<\infty$ & $2$\\
$n\ge8$ & $3\le M(n)<\infty$; exact value open & $3$ or $4$\\
\bottomrule
\end{tabular}
\caption{State-complexity consequences of the exact obstruction thresholds and the finite-horizon existence theorem.  The last path entry is an exact two-value alternative, not an estimate.}
\label{tab:state-complexity}
\end{table}


The unresolved three-state path problem has an exact dual interpretation.  If a three-state controller succeeds on every finite induced path, then $M_{\mathcal P,1}(n)=3$ for every $n\ge8$.  If no such controller exists, then the finite set of three-state tables has a finite maximum first-trap size $F_{\mathcal P,1}(3)$, and Theorem~\ref{thm:FM-duality} implies $M_{\mathcal P,1}(n)=4$ for all sufficiently large $n$.  Thus deciding universal three-state path rendezvous is equivalent to deciding whether the eventual path state complexity is three or four.

\subsection{Finite state complexity at every fixed horizon}

The upper bound in \eqref{eq:M7-bound} is a general finite-horizon fact.  It must not be confused with a single finite controller that solves mazes of all sizes.

\begin{theorem}[Finite-horizon solvability]\label{thm:finite-M}
For every fixed finite $n$ and every $\rho\ge1$,
\[
 \boxed{M_{\mathcal G,\rho}(n)<\infty.}
\]
The controller may depend on $n$ but not on the maze or the starting cell.
\end{theorem}
\begin{proof}
Fix $n$. Store a bounded exploration record: visited coordinates relative to the
start, their observed masks, the current relative coordinate, a DFS stack and a
mode flag. Coordinates are updated internally from chosen compass moves, not read
from the environment. Since every cell is within graph distance $n-1$ of the start,
they lie in $[-(n-1),n-1]^2$. The finite collection of possible records is a finite
state set.

At each cell, move to the first unvisited neighbour in a fixed compass order and
push the current coordinate onto the stack; when none remains, traverse the stack
edge to the parent. Read a newly reached cell's mask on the transition that chooses
its next move. Thus DFS needs no stationary processing step and eventually returns
to the root with the complete relative map.

Choose the lexicographically largest mapped cell and its first open neighbour in
fixed compass order. Maps from different starts differ by translation, which
preserves both choices; both agents therefore select the same physical edge.
Follow a stored route to its selected endpoint, then shuttle on that edge forever.
The mode change takes the first route edge, or the selected edge itself when already
at the endpoint, so every transition moves. Once both agents shuttle, their
Manhattan distance is zero or one, proving success for every $\rho\ge1$.

A singleton succeeds at time zero. Complete impossible-record rows by arbitrary
legal moves. The resulting table depends on $n$, not on the maze or either start.
\end{proof}

The radius restriction in Theorem~\ref{thm:finite-M} is genuine for this model: when $0\le\rho<1$, opposite starts on a two-cell edge swap forever under compulsory synchronous motion, independently of the number of states.

The next two sections fix first a complete test menu and then one host for all
two-state controllers, strengthening the obstruction quantifiers in the same model.
